Fix the class constants and a law \(P\) satisfying the stated conditions.  All constants below may depend on those fixed constants, on \(d,\alpha,\beta,\gamma\), and on \(x_0\), but not on \(P\) or \(n\).

The first assertion is algebraic.  By the definition of the CATE regression,
\[
\mu_{1,P}(x_0)-\mu_{0,P}(x_0)=\tau_P(x_0).
\]
Hence the direct plug-in based on supplied outcome regressions has identically zero error.  By \(\mathrm{lem:cate\mbox{-}identification}\), this observed-data quantity is also the selected continuous version of \(\mathbb E_P[Y(1)-Y(0)\mid X=x_0]\).

We next analyze the separately oracleized sample estimator.  Abbreviate
\[
e(x)=\pi_P(x),\qquad
\eta(x)=e(x)\mu_{1,P}(x)+\{1-e(x)\}\mu_{0,P}(x),
\qquad
\nu(x)=e(x)\{1-e(x)\}.
\]
For \(a\in\{0,1\}\), the definitions of \(\eta\), \(\mu_{a,P}\), and \(\tau_P\) give
\[
\mathbb E_P[Y-\eta(X)\mid A=a,X]
=\{a-e(X)\}\tau_P(X).
\]
Multiplying by \(a-e(X)\), then averaging first over \(Y\) and then over \(A\) conditional on \(X\), yields
\[
\begin{aligned}
\mathbb E_P[(A-e(X))(Y-\eta(X))\mid X]
&=\nu(X)\tau_P(X),\\
\mathbb E_P[\{A-e(X)\}^2\mid X]
&=\nu(X).
\end{aligned}
\]
The overlap assumption gives
\[
\nu(X)\geq\varepsilon(1-\varepsilon)>0.
\]
Consequently the weighted conditional normal equation has the unique regression target \(\tau_P(X)\).  This proves the conditional-target assertion, and \(\mathrm{lem:cate\mbox{-}identification}\) supplies its causal interpretation.

It remains to establish the uniform risk bound.  Put
\[
m=\lceil\gamma\rceil-1,
\]
let \(q\) be the fixed dimension of the polynomials of total degree at most \(m\), abbreviate \(h=h_n\), and write \(\mathcal C=[-1/2,1/2]^d\) and \(b=b_\gamma\).  Since \(x_0\in(0,1)^d\), there is \(h_0>0\) such that
\[
x_0+h\mathcal C\subset[0,1]^d
\quad\text{whenever}\quad 0<h\leq h_0.
\]
The dimensions of the degree-\(\lceil\alpha\rceil-1\) and degree-\(\lceil\beta\rceil-1\) one-cell polynomial spaces are fixed, while \(k_n\to\infty\).  Thus there is a fixed \(n_0\) such that for \(n\geq n_0\) the localization window is interior and none of the finite-sample fallbacks in \(\mathrm{def:rank\mbox{-}capped\mbox{-}estimator}\) is invoked.  In each cyclic fit, replacing both propensity fits by \(e\) and the outcome fit by \(\eta\) removes all nuisance-fold randomness; the corresponding evaluation fold, local-polynomial equation, spectral safeguard, and final truncation remain exactly those in that definition.

Fix one evaluation fold and denote its size by \(N\).  The deterministic balanced split gives \(N\geq n/4\) for every sufficiently large \(n\).  For observations in this fold define
\[
T_i=\frac{X_i-x_0}{h},
\qquad
K_i=h^{-d}\mathbf 1\{T_i\in\mathcal C\}.
\]
The oracle empirical local Gram matrix is
\[
\widehat Q
=\frac1N\sum_i K_i b(T_i)b(T_i)^\top\{A_i-e(X_i)\}^2.
\]
By the two conditional identities above and a change of variables, its expectation is
\[
Q_h
=\int_{\mathcal C}b(t)b(t)^\top
\nu(x_0+ht)f_P(x_0+ht)\,dt.
\]
The coordinates of \(b\) are Lebesgue-orthonormal on \(\mathcal C\).  Therefore the density band and overlap imply, for every \(z\in\mathbb R^q\),
\[
\varepsilon(1-\varepsilon)f_-\lVert z\rVert_2^2
\leq z^\top Q_hz
\leq \frac{f_+}{4}\lVert z\rVert_2^2.
\]
In particular, \(Q_h\) is uniformly nonsingular.  Moreover, its spectrum lies strictly inside the safeguard interval \([\varepsilon^2f_-/2,2f_+]\), because
\[
\varepsilon(1-\varepsilon)f_->\varepsilon^2f_-/2
\quad\text{and}\quad
f_+/4<2f_+.
\]
Choose a constant \(\delta>0\) smaller than both spectral margins and set
\[
G=\{\lVert\widehat Q-Q_h\rVert_{\mathrm{op}}\leq\delta\}.
\]
The fixed-dimensional polynomial vector \(b\) is bounded on the compact cube \(\mathcal C\).  Using \(f_P\leq f_+\), \(|A-e(X)|\leq1\), independence within the fold, and \(\lVert M\rVert_{\mathrm{op}}\leq\lVert M\rVert_{\mathrm F}\), we obtain
\[
\begin{aligned}
\mathbb E_P\lVert\widehat Q-Q_h\rVert_{\mathrm F}^2
&\leq \frac{1}{N}\mathbb E_P
 \left\lVert K_i b(T_i)b(T_i)^\top\{A_i-e(X_i)\}^2\right\rVert_{\mathrm F}^2\\
&\leq \frac{C}{Nh^d}.
\end{aligned}
\]
Hence Markov's inequality gives
\[
P(G^c)\leq\frac{C}{Nh^d}.
\]
On \(G\), Weyl's eigenvalue inequality places the spectrum of \(\widehat Q\) inside the safeguard interval.  Thus the safeguard is inactive and
\[
\mathsf C_{\mathrm{loc}}(\widehat Q)=\widehat Q,
\qquad
\lVert\widehat Q^{-1}\rVert_{\mathrm{op}}\leq C.
\]

Let \(p_h(t)\) be the Taylor polynomial of \(x\mapsto\tau_P(x)\) at \(x_0\), evaluated at \(x=x_0+ht\), through total degree \(m\).  Since \(b\) spans this polynomial space, there is \(\theta_h\in\mathbb R^q\) such that
\[
p_h(t)=b(t)^\top\theta_h,
\qquad
b(0)^\top\theta_h=\tau_P(x_0).
\]
The H\"older CATE assumption gives the uniform Taylor remainder
\[
\sup_{t\in\mathcal C}
|\tau_P(x_0+ht)-p_h(t)|\leq Ch^\gamma.
\]
Define the vector summand
\[
U_i=K_i b(T_i)\left[
(A_i-e(X_i))(Y_i-\eta(X_i))
-\{A_i-e(X_i)\}^2p_h(T_i)
\right].
\]
The conditional-moment identities and the Taylor remainder imply
\[
\begin{aligned}
\lVert\mathbb E_P U_i\rVert_2
&=\left\lVert
\int_{\mathcal C}b(t)\nu(x_0+ht)f_P(x_0+ht)
\{\tau_P(x_0+ht)-p_h(t)\}\,dt
\right\rVert_2\\
&\leq Ch^\gamma.
\end{aligned}
\]
The declared support of \(Y\) gives \(|Y|\leq1\).  Also \(|\eta|\leq1\), and \(|p_h|\leq2+Ch^\gamma\) on \(\mathcal C\).  The boundedness of \(b\) and the density upper bound therefore yield
\[
\mathbb E_P\lVert U_i\rVert_2^2\leq Ch^{-d}.
\]
Independence and Jensen's inequality now give
\[
\mathbb E_P\left\lVert\frac1N\sum_iU_i\right\rVert_2
\leq
\lVert\mathbb E_PU_i\rVert_2
+\left\{
\mathbb E_P\left\lVert
\frac1N\sum_i(U_i-\mathbb E_PU_i)
\right\rVert_2^2
\right\}^{1/2}
\leq Ch^\gamma+C(Nh^d)^{-1/2}.
\]

Let \(\widehat\theta\) be the safeguarded oracle coefficient for this fold.  On \(G\), the safeguard is inactive, so the oracle normal equations give the exact identity
\[
\widehat\theta-\theta_h
=\widehat Q^{-1}\frac1N\sum_iU_i.
\]
Since \(b(0)\) has fixed dimension and norm, it follows that
\[
\mathbb E_P\!\left[
\mathbf 1_G
\left|b(0)^\top\widehat\theta-\tau_P(x_0)\right|
\right]
\leq C\{h^\gamma+(Nh^d)^{-1/2}\}.
\]

Apply this bound to all three evaluation folds and let \(G_{\mathrm{all}}\) be the intersection of their good-Gram events.  Balanced fold sizes and a union bound give
\[
P(G_{\mathrm{all}}^c)\leq C(nh^d)^{-1}.
\]
On \(G_{\mathrm{all}}\), the absolute error of the average of the three fold intercepts is at most the average of their absolute errors.  Projection onto \([-2,2]\) cannot increase distance from \(\tau_P(x_0)\in[-2,2]\).  On \(G_{\mathrm{all}}^c\), the projected estimator and target differ by at most four.  Consequently,
\[
\sup_P
\mathbb E_P\left|
\widehat\tau_n^{\mathrm{or}}(x_0)-\tau_P(x_0)
\right|
\leq
C\left\{h^\gamma+(nh^d)^{-1/2}+(nh^d)^{-1}\right\},
\]
where the supremum is over all laws satisfying the stated conditions.  Finally, \(h=h_n=n^{-1/(2\gamma+d)}\) implies
\[
h^\gamma=(nh^d)^{-1/2}=n^{-\gamma/(2\gamma+d)},
\qquad
(nh^d)^{-1}=n^{-2\gamma/(2\gamma+d)}.
\]
For \(n<n_0\), both the retained truncation and \(\tau_P(x_0)\in[-2,2]\) bound the error by four, so enlargement of the constant covers these finitely many indices.  We conclude uniformly over the stated class that
\[
\mathbb E_P\left|
\widehat\tau_n^{\mathrm{or}}(x_0)-\tau_P(x_0)
\right|
=O\!\left(n^{-\gamma/(2\gamma+d)}\right),
\]
as claimed.